\documentclass[unicode, 12pt, a4paper, oneside, fleqn]{article}

% This file is self-contained so it can also be compiled by editors that
% compile a single .tex file without copying the neighboring project folders.
\usepackage{fontspec}
\usepackage{polyglossia}
\setdefaultlanguage{russian}
\setotherlanguage{english}
\setmainfont{DejaVu Serif}
\setmonofont{DejaVu Sans Mono}
\usepackage{geometry}
\geometry{
  a4paper,
  left=2.54cm,
  right=2.54cm,
  top=2.54cm,
  bottom=2.54cm,
  headheight=0mm,
  headsep=0mm,
  footnotesep=0mm,
  footskip=0mm,
}
\usepackage{nopageno}
\usepackage{titlesec}
\titleformat*{\section}{\large\bfseries}
\titlespacing\section{0pt}{12pt plus 4pt minus 2pt}{0pt plus 2pt minus 2pt}
\setcounter{secnumdepth}{0}
\usepackage{enumitem}
\usepackage{amsmath}
\usepackage{amssymb}
\usepackage{hyperref}
\usepackage{longtable}
\usepackage[table]{xcolor}
\usepackage{array}
\usepackage{listings}
\usepackage{alltt}
\usepackage{csquotes}
\DeclareQuoteStyle{russian}
  {\guillemotleft}{\guillemotright}[0.025em]
  {\quotedblbase}{\textquotedblleft}
\ExecuteQuoteOptions{style=russian}
\usepackage{graphicx}
\usepackage{pgfplots}
\pgfplotsset{compat=1.18}

\lstset{
  tabsize=2,
  breaklines,
  columns=fullflexible,
  flexiblecolumns,
  numbers=left,
  numberstyle={\footnotesize},
  extendedchars,
  inputencoding=utf8,
  escapechar=@,
  basicstyle=\footnotesize\ttfamily,
  showstringspaces=false
}

\parindent=0pt
\parskip=3pt

\begin{document}

\begin{titlepage}
\begin{center}
\large
МОСКОВСКИЙ АВИАЦИОННЫЙ ИНСТИТУТ\\ (НАЦИОНАЛЬНЫЙ ИССЛЕДОВАТЕЛЬСКИЙ УНИВЕРСИТЕТ)

\vspace{20pt}

Институт №8 <<Компьютерные науки и прикладная математика>>

Кафедра 806 <<Вычислительная математика и программирование>>
\end{center}

\vspace{60pt}

\begin{center}
\bfseries
\large
Лабораторная работа № 3 \\по курсу \enquote{Операционные системы}

\vspace{54pt}
\end{center}

\vfill

\begin{flushright}
\large
\begin{tabular}{rl}
Выполнил: & Ягид С. Е. \\
Группа: & М8О-206БВ-25 \\
Проверил: & Ефимов Д. А. \\
\end{tabular}
\end{flushright}

\vspace{92pt}

\begin{center}
\large
Москва, \the\year
\end{center}
\end{titlepage}

\pagebreak

\section{Условие}
{\bfseries Цель работы:} получить практические навыки создания потоков операционной системы, распределения вычислений между потоками и измерения ускорения параллельной программы.

{\bfseries Задание:} вычислить площадь круга методом Монте-Карло. Ограничение числа одновременно работающих потоков задаётся при запуске. В программе выводятся идентификаторы процесса и потоков; эксперимент сравнивает время работы при разном количестве потоков.

{\bfseries Вариант:} 16. Пользователь задаёт радиус круга.

\section{Метод решения}
Для оценки площади равномерно выбираются $N$ точек $(x,y)$ в квадрате $[-1,1]\times[-1,1]$. Точка считается попавшей в круг единичного радиуса, если $x^2+y^2\leq 1$. Если $H$ точек попало внутрь, доля $H/N$ приближает отношение площади круга к площади квадрата. Поэтому площадь круга радиуса $R$ оценивается формулой
\[
    A \approx 4R^2\frac{H}{N}.
\]

Общее число точек делится между $T$ потоками: каждый получает $q=\lfloor N/T\rfloor$ точек, а первые $r=N\bmod T$ потоков получают ещё по одной. Каждый поток хранит собственное число попаданий; общий счётчик во время цикла не изменяется, поэтому блокировка для него не нужна. Главный поток дожидается завершения рабочих через \texttt{WaitForMultipleObjects}, затем складывает локальные результаты.

\section{Описание программы}
Исходный файл программы — \texttt{main.c}. Структура \texttt{WorkerData} хранит размер участка, локальный результат, начальное состояние генератора случайных чисел и TID. Функция \texttt{count\_points} обрабатывает участок одного рабочего потока. Генератор псевдослучайных чисел имеет отдельное состояние у каждого потока.

Для создания и ожидания потоков используются Win32-функции \texttt{CreateThread} и \texttt{WaitForMultipleObjects}; дескрипторы закрываются через \texttt{CloseHandle}. Идентификаторы процесса и потоков получаются через \texttt{GetCurrentProcessId} и \texttt{GetCurrentThreadId}. Время измеряется высокоточным счётчиком \texttt{QueryPerformanceCounter} и частотой \texttt{QueryPerformanceFrequency}. Ошибки системных вызовов проверяются, память для данных рабочих освобождается через \texttt{free}.

Ключевые фрагменты реализации:
\begin{lstlisting}[language=C]
const uint64_t base_count = sample_total / thread_count;
const uint64_t remainder = sample_total % thread_count;
workers[i].sample_count = base_count + (i < remainder ? 1 : 0);
\end{lstlisting}

\begin{lstlisting}[language=C]
if (x * x + y * y <= 1.0) {
    ++hits;
}
data->hits = hits; /* отдельный результат этого потока */
\end{lstlisting}

\section{Результаты}
Программа собрана командой \texttt{gcc -std=c11 -O2 -Wall -Wextra -Wpedantic main.c -o monte\_carlo.exe}. Формат запуска: \texttt{monte\_carlo.exe <радиус> <число\_точек> <число\_потоков>}. Для замеров радиус равен 1, число точек — 20,000,000; каждое измерение повторялось три раза. Время включает создание потоков и ожидание их завершения. В таблице указана медиана трёх измерений.

\begin{center}
\begin{tabular}{r r r r r}
\hline
Потоки $p$ & Замеры, мс & Медиана $T_p$, мс & Ускорение $S_p$ & Эффективность $E_p$ \\
\hline
1 & 73.681; 70.477; 70.666 & 70.666 & 1.000 & 1.000 \\
2 & 35.379; 38.743; 40.121 & 38.743 & 1.824 & 0.912 \\
4 & 19.303; 18.838; 19.597 & 19.303 & 3.661 & 0.915 \\
8 & 18.016; 19.351; 18.343 & 18.343 & 3.852 & 0.482 \\
\hline
\end{tabular}
\end{center}

Ускорение вычислено как $S_p=T_1/T_p$, эффективность — как $E_p=S_p/p$.

\begin{center}
\begin{tikzpicture}
\begin{axis}[
  width=0.85\textwidth,
  height=6cm,
  xlabel={Число потоков},
  ylabel={Время, мс},
  xtick={1,2,4,8},
  grid=both,
  ymin=0,
  xmin=0,
  xmax=9,
  mark size=2.5pt
]
\addplot[blue, thick, mark=*] coordinates {
  (1,70.666) (2,38.743) (4,19.303) (8,18.343)
};
\end{axis}
\end{tikzpicture}
\end{center}

При четырёх потоках время стало примерно в 3.66 раза меньше, чем при одном. Переход от четырёх к восьми потокам дал небольшое дополнительное ускорение; эффективность при восьми потоках снизилась. Это может быть связано с ограниченным числом доступных вычислительных ядер и накладными расходами планирования. Число ядер в этой среде получить не удалось, поэтому причина формулируется как предположение. Во всех запусках оценка площади круга радиуса 1 находилась близко к $\pi$; разброс между запусками ожидаем из-за случайной выборки.

Пример запуска на 1,000,000 точек с четырьмя потоками: оценка площади $3.140772$, время $1.091$ мс. При запуске программа выводит PID процесса и TID каждого рабочего потока. Эти идентификаторы позволяют сопоставить выполнение потоков с событиями 	exttt{Thread Create} и 	exttt{Thread Exit} в Process Monitor.

\section{Выводы}
Реализована параллельная оценка площади круга методом Монте-Карло. Работа делится между потоками без совместного изменения счётчика попаданий; результаты объединяются после ожидания потоков. На данном измерении увеличение числа потоков с одного до четырёх заметно сократило время расчёта, а восемь потоков дали существенно меньший дополнительный выигрыш. Оценка площади приближается к точному значению при увеличении числа случайных точек.

\pagebreak

\section{Исходная программа}
Полный исходный код находится в файле \texttt{main.c}. В основной текст включены фрагменты, важные для распределения работы и подсчёта попаданий.

\end{document}
